\chapter{Parallel Repetition: Driving the Soundness Down}
\label{lec:raz}

{\sffamily\large\color{softgray}Playing a game $k$ times at once is not the same as playing it $k$ times\par}
\vspace{0.5em}

\providecommand{\lamp}{\lambda_{+}}

\noindent
Lecture~\ref{lec:why} ended with a promise.  The PCP theorem gives a game of value $1$ or $1-\eta_0$; repeating it $k=O(\log\frac1\delta)$ times in parallel gives a game of value $1$ or at most $\delta$,
at the price of an alphabet of size $\poly(1/\delta)$.  The proof there was two lines long: ``apply Theorem~\ref{thm:raz}''.  This lecture opens that black box.

The difficulty is easy to state.  A verifier that repeats itself \emph{sequentially} multiplies its acceptance probabilities (Proposition~\ref{prop:pcp-equiv}).
Two provers cannot be questioned sequentially $k$ times without $k$ rounds of interaction, so the referee sends all $k$ questions at once, and each prover sees all $k$ of its questions before answering any of them.
It may therefore answer coordinate~$1$ using what it learned about coordinate~$2$.  Does that help the provers?  It does (Section~\ref{sec:notpower}), and for years it was open whether the value even tends to $0$; a non-constructive proof that it does came first, as Raz notes in his introduction.
\citet{Raz98} proved that it tends to $0$ exponentially fast in $k$, with a proof of about forty pages.

We shall do three things.  First, we meet the examples that explain why the answer is not obvious: a two-answer game whose two-fold repetition beats the square of its value, and a game of value $\frac12$ whose two-fold repetition still has value~$\frac12$.
Second, we prove a repetition theorem \emph{completely}, for projection games whose ``collision graph'' is an expander.  The proof, due to \citet{DinurSteurer14}, is linear algebra: it takes a few pages, and it is enough to establish the Label Cover theorem
(Theorem~\ref{thm:lc}) with explicit constants, using a one-line trick that makes any game expanding.  Third, we state what is known without the expansion hypothesis, sketch how Raz and Holenstein prove it, and explain why repetition can never,
by itself, prove the Unique Games Conjecture.

\begin{goalbox}{After this lecture you should be able to\ldots}
\begin{enumerate}[leftmargin=*,label=(\arabic*)]
  \item compute the $k$-fold value of a small game by exhaustive search, and exhibit a strategy that beats $\val(G)^k$ (Section~\ref{sec:notpower});
  \item define the collision value $\|G\|$ and prove $\val(G)^2\le\|G\|^2\le\val(G)$ (Section~\ref{sec:operator});
  \item prove that $\|G\otimes H\|\le\|G\|$ and, more sharply, $\|G\otimes H\|\le\lamp(G)\,\|H\|$ (Sections~\ref{sec:tensor}--\ref{sec:product});
  \item prove that $\lamp(G)\le1-\gamma\eta/54$ for a $\gamma$-expanding projection game of value $\le1-\eta$, and deduce the Label Cover theorem with an explicit number of repetitions (Sections~\ref{sec:expander}--\ref{sec:lc});
  \item state the repetition theorems of Raz, Holenstein, Rao and Dinur--Steurer, say what fails (strong repetition) and why that obstructs a repetition-based proof of the Unique Games Conjecture (Sections~\ref{sec:general}--\ref{sec:ugc-rep}).
\end{enumerate}
\end{goalbox}


%% ===========================================================================
\section{Repeating a game is not raising its value to a power}
\label{sec:notpower}
%% ===========================================================================

We use the games of Lecture~\ref{lec:why} (Definition~\ref{def:game}) with one harmless generalization: the referee draws the edge from an arbitrary probability distribution~$\mu$ rather than the uniform one.
All games in Lectures~\ref{lec:why}--\ref{lec:games} have uniform~$\mu$, and a rational~$\mu$ can always be made uniform by repeating edges (see Remark~\ref{rem:multiplicity} below).

\begin{definition}[Weighted projection game]\label{def:wgame}
A \defn{projection game} $G$ has question sets $U$ (Alice) and $W$ (Bob), label sets $\Sigma_U,\Sigma_W$, a finite set $E$ of edges $e=(u_e,w_e)\in U\times W$ (parallel edges allowed) with a probability distribution $\mu$ on $E$, and for every edge a map $\pi_e\colon\Sigma_U\to\Sigma_W$.
The referee draws $e\sim\mu$, sends $u_e$ to Alice and $w_e$ to Bob, and accepts answers $a\in\Sigma_U$, $b\in\Sigma_W$ iff $\pi_e(a)=b$.
We write $\mu_U,\mu_W$ for the marginals of $\mu$ on questions (assumed positive), $\mu(\cdot\mid w)$ for $\mu$ conditioned on $w_e=w$, and $\val(G)$ for the maximum acceptance probability over deterministic strategies.
The \defn{product} $G\otimes H$ has questions $U\times U'$ and $W\times W'$, labels $\Sigma_U\times\Sigma_{U'}$ and $\Sigma_W\times\Sigma_{W'}$, edges $E\times E'$ with the product distribution, and constraint maps $\pi_e\times\pi_{e'}$; $G^{\otimes k}$ is the $k$-fold product.
\end{definition}

\noindent
This is the $k$-fold parallel repetition of Lecture~\ref{lec:why}: the referee draws $k$ independent edges, sends all $k$ Alice-questions to Alice and all $k$ Bob-questions to Bob, and accepts iff every coordinate is accepted.
Note that \emph{totality} of $\pi_e$ is part of the definition: every label of Alice has exactly one accepted partner.  Unique Games are the case of bijective maps.

\begin{remark}\label{rem:multiplicity}
If $\mu$ takes rational values, write $\mu(e)=n_e/N$; replacing $e$ by $n_e$ parallel copies gives an equivalent game with the uniform distribution on $N$ edges.
\end{remark}

\subsection{The easy inequalities}

\begin{lemma}[Elementary facts]\label{lem:rep-basic}
For every game $G$ and $k\ge2$: $\ \val(G)^k\ \le\ \val(G^{\otimes k})\ \le\ \val(G^{\otimes(k-1)})$.  In particular $\val(G)=1$ implies $\val(G^{\otimes k})=1$.
\end{lemma}
\begin{proof}
Lower bound: both provers play an optimal strategy of $G$ independently in each coordinate; the $k$ coordinates are independent.
Upper bound: take a deterministic strategy $(f,g)$ for $G^{\otimes k}$ and write $e=(e_1,\dots,e_k)$.  Winning all $k$ coordinates implies winning the first $k-1$, so
$\val(G^{\otimes k})\le\E_{e_k}\Pr_{e_1..e_{k-1}}[\text{first }k-1\text{ coordinates won}\mid e_k]$, hence some fixed edge $e_k=(u_k,w_k)$ attains at least this average.  With $u_k$ and $w_k$ hard-wired, $f(\cdot,u_k)$ and $g(\cdot,w_k)$ are strategies for $G^{\otimes(k-1)}$ that win with probability at least $\val(G^{\otimes k})$.
\end{proof}

\noindent
Sequential repetition would give $\val(G)^k$ exactly.  Is the first inequality an equality?  Here is the smallest example showing that it is not.

\begin{example}[CHSH]\label{ex:chsh}
Let $U=W=\{0,1\}$ and $\Sigma_U=\Sigma_W=\{0,1\}$, with the uniform distribution on the four pairs $(x,y)$ and the constraint $a\oplus b=x\wedge y$, \ie $\pi_{(x,y)}(a)=a\oplus xy$.  Every $\pi_{(x,y)}$ is a bijection, so this is a Unique Game with two labels.

\emph{Its value is $\frac34$.}  The four constraints $f(x)\oplus g(y)=xy$ cannot all hold: adding them, the left sides cancel in pairs and give $0$, while the right sides add up to $0+0+0+1=1$.  The strategy $f=g=0$ satisfies three.

\emph{Its two-fold repetition has value $\frac58>\big(\frac34\big)^2=\frac9{16}$.}  Questions are now $x=(x_1,x_2)$ and $y=(y_1,y_2)$, answers are pairs, and both coordinates must satisfy $a_i\oplus b_i=x_iy_i$.  Consider the strategy
\[
  a(x)=\begin{cases}(1,0)&x=(1,1)\\(0,0)&\text{otherwise}\end{cases}\qquad
  b(y)=\begin{cases}(0,1)&y=(1,1)\\(0,0)&\text{otherwise.}\end{cases}
\]
Table~\ref{tab:chsh-win} shows that it wins on $10$ of the $16$ question pairs: $10/16=\frac58$.
Exhaustive search over all $4^4\cdot4^4$ strategies confirms that nothing does better; for $k=3$ an exhaustive search over Alice's $8^8$ strategies, with Bob responding optimally, gives $\val(G^{\otimes3})=\frac{31}{64}$, which exceeds \mbox{$\big(\frac34\big)^3=\frac{27}{64}$}.
\end{example}

\begin{table}[t]
\centering\small
\renewcommand{\arraystretch}{1.15}
\begin{tabular}{@{}c|cccc@{}}
\toprule
$x\backslash y$ & $00$ & $01$ & $10$ & $11$\\
\midrule
$00$ & \checkmark & \checkmark & \checkmark & $\cdot$\\
$01$ & \checkmark & $\cdot$ & \checkmark & \checkmark\\
$10$ & \checkmark & \checkmark & $\cdot$ & $\cdot$\\
$11$ & $\cdot$ & $\cdot$ & \checkmark & \checkmark\\
\bottomrule
\end{tabular}
\caption{The winning question pairs of the two-fold CHSH strategy of Example~\ref{ex:chsh}: $10$ of $16$.  A single coordinate of the same strategy would win only $\frac34$ of the time on its own; the strategy gains by letting each prover's answer in coordinate~$i$ depend on the \emph{other} coordinate's question.}
\label{tab:chsh-win}
\end{table}

\noindent
The gain is real but small: $\frac58$ against $\frac9{16}$.  The next example shows that without structure, repetition can fail to help at all.

\begin{example}[Feige's non-interactive agreement game]\label{ex:feige}
Let $x,y\in\{0,1\}$ be independent uniform questions to Alice and Bob, and let the label set of both be $\{A_0,A_1,B_0,B_1\}$.  The provers win iff they give the \emph{same} label and it is either $A_x$ (``we agree on Alice's question'') or $B_y$ (``we agree on Bob's question'').

\emph{Its value is $\frac12$.}  Fix deterministic strategies $f,g$.  For a given Bob-question $y$, the label $g(y)$ can win against at most one $x$ if $g(y)=A_s$ (namely $x=s$), against at most two if $g(y)=B_y$ (namely those $x$ with $f(x)=B_y$), and against none otherwise.
If some $y$ is won twice then $g(y)=B_y$ and $f\equiv B_y$, and then the other $y'$ is never won; otherwise each $y$ is won at most once.  In either case at most $2$ of the $4$ question pairs are won.  Always answering $A_0$ achieves $\frac12$.

\emph{Its two-fold repetition has value $\frac12$ as well.}  Let Alice, on $(x_1,x_2)$, answer $(A_{x_1},B_{x_1})$, and Bob, on $(y_1,y_2)$, answer $(A_{y_2},B_{y_2})$.  In coordinate~$1$ the pair $(A_{x_1},A_{y_2})$ is accepted iff $y_2=x_1$ (the label $A_{x_1}$ is Alice's own question); in coordinate~$2$ the pair $(B_{x_1},B_{y_2})$ is accepted iff $y_2=x_1$ (the label $B_{y_2}$ is Bob's own question).  Both conditions are the same event $\{x_1=y_2\}$, which has probability $\frac12$.
Together with Lemma~\ref{lem:rep-basic}, $\val(G^{\otimes2})=\val(G)=\frac12$.  An exhaustive search over all strategies confirms both numbers.
\end{example}

\begin{pitfall}{repetition does not have to help, and does not help as the product says}
Example~\ref{ex:feige} is a game for which \emph{no} bound of the form $\val(G^{\otimes2})\le\varphi(\val(G))$ with $\varphi(x)<x$ can hold; a repetition theorem must therefore speak about $k$ large, and about the size of the answer sets or the structure of the constraints.
Example~\ref{ex:chsh} shows that even for Unique Games the true value of the repeated game lies strictly between $\val(G)^k$ and $\val(G)$.  Neither game is exotic: the first is a free game on two questions, the second is the XOR game behind Bell's inequality.
The Feige game is \emph{not} a projection game in our sense (a label $A_{1-x}$ of Alice has no accepted partner at all), and Lemma~\ref{lem:val-coll} below uses totality; this is no accident.
\end{pitfall}

We will return to these two examples in Section~\ref{sec:general}.  For now the lesson is: to prove that $\val(G^{\otimes k})\to0$ we cannot reason one coordinate at a time.  We need a quantity of $G$ that \emph{is} multiplicative-like under products and still controls the value.


%% ===========================================================================
\section{A projection game as a linear operator}
\label{sec:operator}
%% ===========================================================================

\subsection{Strategies, the operator $P_G$, and the success probability}

Let $G$ be a projection game.  An \defn{assignment} for Alice is a function $f\colon U\times\Sigma_U\to[0,1]$ with $\sum_af(u,a)=1$ for every $u$; $f(u,a)$ is the probability that Alice answers $a$ to $u$.
A deterministic \defn{labeling} $\ell\colon U\to\Sigma_U$ is the assignment $f_\ell(u,a)=\mathbf 1[a=\ell(u)]$.  Assignments for Bob are defined in the same way on $W\times\Sigma_W$.
On real functions $f,f'$ on $U\times\Sigma_U$ and $g,g'$ on $W\times\Sigma_W$ we use the inner products
\[
  \langle f,f'\rangle=\E_{u\sim\mu_U}\sum_af(u,a)f'(u,a),\qquad
  \langle g,g'\rangle=\E_{w\sim\mu_W}\sum_bg(w,b)g'(w,b),
\]
and the norms $\|f\|=\langle f,f\rangle^{1/2}$.
The \defn{operator of the game} sends a function on $U\times\Sigma_U$ to a function on $W\times\Sigma_W$:
\[
  (P_Gf)(w,b)\;=\;\E_{e\sim\mu(\cdot\mid w)}\ \sum_{a\colon\pi_e(a)=b}f(u_e,a).
\]
In words: to evaluate $P_Gf$ at $w$, pick a random neighbour of $w$, let Alice answer according to $f$, and project her answer to a label of Bob.  Because every $\pi_e$ is a \emph{total} map, $\sum_b(P_Gf)(w,b)=\E_{e\mid w}\sum_af(u_e,a)=1$ when $f$ is an assignment:
\emph{$P_Gf$ is again an assignment}, this time for Bob.

\begin{lemma}[The success probability is a bilinear form]\label{lem:bilinear}
If Alice plays the assignment $f$ and Bob plays $g$, the referee accepts with probability $\langle g,P_Gf\rangle$.  Consequently $\val(G)=\max\langle g,P_Gf\rangle$ over assignments, and the maximum is attained by deterministic ones.
\end{lemma}
\begin{proof}
The acceptance probability is $\E_{e\sim\mu}\sum_{a,b}f(u_e,a)g(w_e,b)\mathbf 1[\pi_e(a)=b]$.  Condition on $w_e=w$ to get $\E_w\sum_bg(w,b)\,\E_{e\mid w}\sum_{a\colon\pi_e(a)=b}f(u_e,a)=\langle g,P_Gf\rangle$.
This expression is bilinear in $(f,g)$, and the set of assignments is a product of simplices whose extreme points are the deterministic ones.
\end{proof}

\subsection{The collision value}

The optimal $g$ for a given $f$ is obvious (answer the most likely label of $P_Gf$), so the interesting object is $P_Gf$ itself.  Its length is the quantity that will behave well under products.

\begin{definition}[Collision value]\label{def:coll}
$\|G\|\;=\;\max\{\|P_Gf\|:\ f\text{ an assignment for Alice}\}$.
\end{definition}

\noindent
Since $f\mapsto\|P_Gf\|$ is convex, the maximum is attained at a deterministic labeling.  The name comes from the following experiment, in which Bob has been eliminated.

\begin{lemma}[The collision game]\label{lem:collgame}
$\|P_Gf\|^2$ is the probability of acceptance in this experiment: draw $w\sim\mu_W$ and two edges $e,e'\sim\mu(\cdot\mid w)$ independently; let two independent copies of Alice answer $a\sim f(u_e,\cdot)$ and $a'\sim f(u_{e'},\cdot)$; accept iff $\pi_e(a)=\pi_{e'}(a')$.
\end{lemma}
\begin{proof}
Expand the square and exchange the sum over $b$ with the sums over $a,a'$:
\begin{align*}
  \|P_Gf\|^2&=\E_w\sum_b\Big(\E_{e\mid w}\sum_{a\colon\pi_e(a)=b}f(u_e,a)\Big)^2\\
  &=\E_w\ \E_{e,e'\mid w}\ \sum_{a,a'}f(u_e,a)\,f(u_{e'},a')\,\mathbf 1\big[\pi_e(a)=\pi_{e'}(a')\big].\qedhere
\end{align*}
\end{proof}

\noindent
For the clause--variable game of a $3$-CNF this is the following experiment (Figure~\ref{fig:collision}): choose a random variable $x$ and two random clauses containing it; the two clause-provers must give satisfying assignments of their clauses that \emph{agree on $x$}.
The same symmetrization appears, implicitly, throughout the PCP literature: for example when a ``line versus point'' low-degree test is turned into a ``line versus line'' test \citep[footnote to Section~2.3]{DinurSteurer14}.

\begin{figure}[t]
\centering
\begin{tikzpicture}[x=1cm,y=1cm,every node/.style={font=\small},
   cl/.style={draw,rounded corners=2pt,minimum width=5.3cm,minimum height=0.95cm,fill=accent!8,align=center},
   vr/.style={draw,circle,minimum size=0.9cm,fill=warn!10},
   ok/.style={font=\bfseries\small,color=green!50!black},
   no/.style={font=\bfseries\small,color=claim}]
  \node[font=\sffamily\small\bfseries,color=accent,anchor=west] at (-2.7,2.55) {(a)\ \ a collision};
  \node[cl] (c1) at (0,1.5) {$C_1=(x_1\vee x_2\vee\neg x_3)$\\ label $(x_1,x_2,x_3)=(1,0,0)$};
  \node[cl] (c2) at (0,0) {$C_2=(\neg x_1\vee x_3\vee x_4)$\\ label $(x_1,x_3,x_4)=(0,0,0)$};
  \node[vr] (x3) at (5.2,0.75) {$x_3$};
  \draw[-{Stealth[length=2.2mm]},thick] (c1.east) -- node[above,sloped,font=\footnotesize] {$x_3=0$} (x3);
  \draw[-{Stealth[length=2.2mm]},thick] (c2.east) -- node[below,sloped,font=\footnotesize] {$x_3=0$} (x3);
  \node[ok,anchor=west] at (5.9,0.75) {agree: accept};
  \begin{scope}[yshift=-3.9cm]
  \node[font=\sffamily\small\bfseries,color=accent,anchor=west] at (-2.7,2.55) {(b)\ \ no collision};
  \node[cl] (d1) at (0,1.5) {$C_1=(x_1\vee x_2\vee\neg x_3)$\\ label $(x_1,x_2,x_3)=(1,0,0)$};
  \node[cl] (d2) at (0,0) {$C_2=(\neg x_1\vee x_3\vee x_4)$\\ label $(x_1,x_3,x_4)=(1,1,0)$};
  \node[vr] (y3) at (5.2,0.75) {$x_3$};
  \draw[-{Stealth[length=2.2mm]},thick] (d1.east) -- node[above,sloped,font=\footnotesize] {$x_3=0$} (y3);
  \draw[-{Stealth[length=2.2mm]},thick] (d2.east) -- node[below,sloped,font=\footnotesize] {$x_3=1$} (y3);
  \node[no,anchor=west] at (5.9,0.75) {disagree: reject};
  \end{scope}
\end{tikzpicture}
\caption{The collision game of the clause--variable game of Figure~\ref{fig:cv}.  The referee picks the variable $x_3$ and two clauses containing it; the two copies of Alice answer with satisfying assignments of their clauses (both are satisfying in (a) and (b)) and are accepted iff their projections to $x_3$ coincide.  Bob does not appear.}
\label{fig:collision}
\end{figure}

\begin{lemma}[Collision value versus value]\label{lem:val-coll}
$\ \val(G)^2\ \le\ \|G\|^2\ \le\ \val(G).$
\end{lemma}
\begin{proof}
Let $f,g$ be deterministic and optimal, so $\val(G)=\langle g,P_Gf\rangle\le\|g\|\,\|P_Gf\|\le\|G\|$, because $\|g\|^2=\E_w\sum_bg(w,b)^2=1$.
For the other inequality let $f$ attain $\|G\|$.  Since $P_Gf$ is an assignment for Bob, Lemma~\ref{lem:bilinear} gives $\|G\|^2=\langle P_Gf,P_Gf\rangle\le\val(G)$.
\end{proof}

\noindent
In particular $\val(G)=1$ iff $\|G\|=1$.  The squaring is harmless for us: ``exponentially small'' survives a square root, at the cost of a factor $2$ in the number of repetitions.
(The proof of the second inequality is the place where totality of the maps $\pi_e$ is used; for a partial projection one must increase some entries of $P_Gf$ to turn it into an assignment, and this is how \citet{DinurSteurer14} proceed.)

\begin{example}[CHSH, by hand]\label{ex:chsh-coll}
In the collision game of Example~\ref{ex:chsh} the referee picks a uniform $y$ and two independent uniform $x,x'$; a deterministic Alice answers $f(x),f(x')$; the pair is accepted iff $f(x)\oplus xy=f(x')\oplus x'y$.
If $x=x'$ (probability $\frac12$) this always holds.  If $x\ne x'$ it holds iff $f(0)\oplus f(1)=y$, which happens for exactly one of the two values of $y$.  So $\|P_Gf\|^2=\frac12+\frac12\cdot\frac12=\frac34$ for every deterministic $f$, \ie
$\|G\|^2=\frac34=\val(G)$.  In Lemma~\ref{lem:val-coll} the left inequality is strict here ($\frac9{16}<\frac34$) and the right one is an equality.
\end{example}


%% ===========================================================================
\section{The collision value under products}
\label{sec:tensor}
%% ===========================================================================

For two functions on question--label pairs the product game acts coordinatewise.  Write $L(X)$ for the real functions on $X$ with the inner product above.

\begin{lemma}[Products are tensor products]\label{lem:tensor}
$P_{G\otimes H}=P_G\otimes P_H$: for $F\colon(U\times U')\times(\Sigma_U\times\Sigma_{U'})\to\R$,
\[
  (P_{G\otimes H}F)\big((w,w'),(b,b')\big)=\E_{e\sim\mu(\cdot\mid w)}\ \E_{e'\sim\mu'(\cdot\mid w')}\ \sum_{a\colon\pi_e(a)=b}\ \sum_{a'\colon\pi_{e'}(a')=b'}F\big((u_e,u_{e'}),(a,a')\big).
\]
\end{lemma}
\begin{proof}
Given $(w,w')$, the edge of $G\otimes H$ is a product of an edge of $G$ and an edge of $H$, drawn independently from $\mu(\cdot\mid w)$ and $\mu'(\cdot\mid w')$.
\end{proof}

\noindent
For the value this tensor structure is of little use, since the value of a product is not the product of the values (Section~\ref{sec:notpower}).  For the \emph{collision} value we can at least prove the monotonicity that is trivial for the value.

\begin{lemma}[Playing in parallel cannot raise the collision value]\label{lem:coll-mono}
For all projection games $G,H$: $\ \|G\otimes H\|\le\|G\|$ and $\|G\otimes H\|\le\|H\|$.
\end{lemma}
\begin{proof}
By symmetry it suffices to prove the first.  Let $F$ be an assignment for $G\otimes H$.  For $u'\in U'$ and $w'\in W'$ define functions on $U\times\Sigma_U$:
\[
  F_{u'}(u,a)=\sum_{a'}F\big((u,u'),(a,a')\big),\qquad F_{w'}=\E_{e'\sim\mu'(\cdot\mid w')}F_{u_{e'}}.
\]
$F_{u'}$ is an assignment for $G$ (non-negative, and $\sum_aF_{u'}(u,a)=\sum_{a,a'}F=1$), hence so is the average $F_{w'}$.
Fix $(w,w')$ and $b$, and put $x_{b'}=(P_{G\otimes H}F)((w,w'),(b,b'))\ge0$ as in Lemma~\ref{lem:tensor}.  Summing over $b'$ removes the constraint on $a'$:
\[
  \sum_{b'}x_{b'}=\E_{e\mid w}\ \E_{e'\mid w'}\sum_{a\colon\pi_e(a)=b}F_{u_{e'}}(u_e,a)=\E_{e\mid w}\sum_{a\colon\pi_e(a)=b}F_{w'}(u_e,a)=(P_GF_{w'})(w,b).
\]
For non-negative numbers $\sum_{b'}x_{b'}^2\le\big(\sum_{b'}x_{b'}\big)^2$.  Therefore
\[
  \|P_{G\otimes H}F\|^2=\E_{w,w'}\sum_{b,b'}x_{b'}^2\ \le\ \E_{w'}\ \E_w\sum_b\big((P_GF_{w'})(w,b)\big)^2=\E_{w'}\|P_GF_{w'}\|^2\ \le\ \|G\|^2.\qedhere
\]
\end{proof}

\noindent
Lemma~\ref{lem:coll-mono} says that $\|G^{\otimes k}\|$ is non-increasing in $k$ (compare Lemma~\ref{lem:rep-basic}).  But it does not say that it \emph{decreases}.  For that we need to compare $G\otimes H$ with $H$ alone, and this is where the argument becomes interesting.


%% ===========================================================================
\section{The product theorem}
\label{sec:product}
%% ===========================================================================

Fix $G$.  We look for a number $\rho_G<1$, depending only on $G$, such that $\|G\otimes H\|\le\rho_G\|H\|$ \emph{for every} projection game~$H$.
Given such a $\rho_G$, $\|G^{\otimes k}\|\le\rho_G^{k-1}\|G\|$ follows at once, and the value of $G^{\otimes k}$ is bounded by Lemma~\ref{lem:val-coll}.

Suppose $F$ is a best assignment for $G\otimes H$ and consider first the ``null'' version of $G$ that never rejects.  This is the trivial game below; comparing $G\otimes H$ with it, the part coming from $H$ will cancel.

\begin{definition}[The trivial game and the parameter $\lamp(G)$]\label{def:lamp}
The \defn{trivial game} $T_G$ has the same Alice side as $G$ (questions $U$, labels $\Sigma_U$, marginal $\mu_U$), Bob's questions are again $U$, Bob has the single label $\ast$, and the edges are $(u,u)$ with probability $\mu_U(u)$.  Every strategy wins.  Its operator computes the mass $(P_{T_G}h)(u,\ast)=\sum_ah(u,a)$ of a function $h\ge0$, so $\|P_{T_G}h\|^2=\E_u\big(\sum_ah(u,a)\big)^2$.
Define
\[
  \lamp(G)\;=\;\sup\Big\{\frac{\|P_Gh\|}{\|P_{T_G}h\|}\ :\ h\colon U\times\Sigma_U\to\R_{\ge0},\ h\ne0\Big\}.
\]
\end{definition}

\noindent
The supremum is over all non-negative functions, not only assignments, so $\|G\|\le\lamp(G)$ (take $h$ an assignment: $\|P_{T_G}h\|=1$).  Also $\lamp(G)\le1$, because for $m(u)=\sum_ah(u,a)$,
$\|P_Gh\|^2=\E_w\sum_b\big(\E_{e\mid w}\sum_{a\colon\pi_e(a)=b}h(u_e,a)\big)^2\le\E_w\big(\E_{e\mid w}m(u_e)\big)^2\le\E_um(u)^2$ by $\sum x_b^2\le(\sum x_b)^2$ and Jensen.

\begin{theorem}[Product theorem; Dinur--Steurer]\label{thm:product}
For all projection games $G,H$: $\ \|G\otimes H\|\le\lamp(G)\cdot\|H\|$.
\end{theorem}
\begin{proof}
Let $F$ be an assignment for $G\otimes H$ with $\|P_{G\otimes H}F\|=\|G\otimes H\|$.  Define a non-negative function $h$ of $x=(u,a)$ and $j=(w',b')$ by applying only the $H$-half of the operator:
\[
  h(x;j)\;=\;\E_{e'\sim\mu'(\cdot\mid w')}\ \sum_{a'\colon\pi_{e'}(a')=b'}F\big((u,u_{e'}),(a,a')\big).
\]
By Lemma~\ref{lem:tensor}, $P_{G\otimes H}F=(P_G\otimes\mathrm{Id})h$ and $P_{T_G\otimes H}F=(P_{T_G}\otimes\mathrm{Id})h$.
The assignment $F$ has the same Alice side in $G\otimes H$ and in $T_G\otimes H$, so it is an assignment for $T_G\otimes H$ as well, and by Lemma~\ref{lem:coll-mono}
\[
  \|(P_{T_G}\otimes\mathrm{Id})h\|=\|P_{T_G\otimes H}F\|\ \le\ \|T_G\otimes H\|\ \le\ \|H\|.
\]
(The left side is positive: for an assignment $F$ the entries of $(P_{T_G}\otimes\mathrm{Id})h$ at $(u,w')$ sum over $b'$ to $1$.)  Dividing,
\[
  \frac{\|G\otimes H\|}{\|H\|}\ \le\ \frac{\|(P_G\otimes\mathrm{Id})h\|}{\|(P_{T_G}\otimes\mathrm{Id})h\|}.
\]
Both numerator and denominator are sums over the columns $j$ of $h$: with $h_j=h(\cdot\,;j)\ge0$ and weights $m_j=\mu_{W'}(w')$,
$\|(P_G\otimes\mathrm{Id})h\|^2=\sum_jm_j\|P_Gh_j\|^2$ and $\|(P_{T_G}\otimes\mathrm{Id})h\|^2=\sum_jm_j\|P_{T_G}h_j\|^2$.
A ratio of sums with positive denominators is at most the largest ratio of corresponding terms, and for each $j$ the ratio is at most $\lamp(G)^2$ by definition (columns with $h_j=0$ contribute $0$ to both sums).
\end{proof}

\noindent
The mechanism is worth a sentence.  The function $h$ summarizes everything that $H$ did to the strategy; the factorization $P_{G\otimes H}=(P_G\otimes\mathrm{Id})(\mathrm{Id}\otimes P_H)$ lets us discard $H$ and keep only the best $h$.
The denominator, supplied by the trivial game, is what turns an inequality about two interacting games into one about~$G$ alone.

\begin{corollary}[Repetition bound in terms of $\lamp$]\label{cor:lamp}
For every projection game $G$ and $k\ge1$: $\ \val(G^{\otimes k})\ \le\ \|G^{\otimes k}\|\ \le\ \lamp(G)^k$.
\end{corollary}
\begin{proof}
Apply Theorem~\ref{thm:product} with $H=G^{\otimes(k-1)}$ repeatedly, then use $\|G\|\le\lamp(G)$ and Lemma~\ref{lem:val-coll}.
\end{proof}

\begin{example}[CHSH, continued]\label{ex:chsh-lamp}
Let $h=m\cdot f_\ell$ for a labeling $\ell$ of Alice's two questions and weights $m=(m_0,m_1)\ge0$.  The computation of Example~\ref{ex:chsh-coll}, now with weights, gives $\|P_Gh\|^2=\frac14(m_0^2+m_1^2)+\frac14m_0m_1$ for \emph{every} $\ell$ (the two questions collide for exactly one of the two values of $y$), while $\|P_{T_G}h\|^2=\frac12(m_0^2+m_1^2)$.
The largest value of the ratio of these two quadratic forms is $\frac34$, attained at $m=(1,1)$.  Deterministic $h$ suffice for the supremum (Step~1 of the proof of Theorem~\ref{thm:expander} below), so $\lamp(G)^2=\frac34$.  Corollary~\ref{cor:lamp} reads $\val(G^{\otimes k})\le(\frac34)^{k/2}$.

This bound is tight for the \emph{collision value}, for every $k$: the product assignment $f^{\otimes k}$ has $P_{G^{\otimes k}}f^{\otimes k}=(P_Gf)^{\otimes k}$, whose norm is $\|P_Gf\|^k=(\sqrt3/2)^k$, so $\|G^{\otimes k}\|=(\sqrt3/2)^k=\lamp(G)^k$.
It is not tight for the value: Table~\ref{tab:chsh-sandwich} puts the chain $\val(G)^k\le\val(G^{\otimes k})\le\|G^{\otimes k}\|=\lamp(G)^k$ next to the true values.
\end{example}

\begin{table}[t]
\centering\small
\renewcommand{\arraystretch}{1.2}
\begin{tabular}{@{}rcccc@{}}
\toprule
$k$ & $\val(G)^k$ & $\val(G^{\otimes k})$ & $\|G^{\otimes k}\|$ & $\lamp(G)^k$\\
\midrule
$1$ & $0.7500$ & $0.7500$ & $0.8660$ & $0.8660$\\
$2$ & $0.5625$ & $0.6250$ & $0.7500$ & $0.7500$\\
$3$ & $0.4219$ & $0.4844$ & $0.6495$ & $0.6495$\\
\bottomrule
\end{tabular}
\caption{CHSH: the true values of Example~\ref{ex:chsh} (second column; exhaustive search) are sandwiched between the na\"ive power (first column) and the collision value (third column), which equals $(\sqrt3/2)^k=\lamp(G)^k$ by Example~\ref{ex:chsh-lamp}.  The first inequality of Lemma~\ref{lem:rep-basic} is strict for $k\ge2$.}
\label{tab:chsh-sandwich}
\end{table}

\begin{pitfall}{$\lamp(G)$ can equal $1$ for a game that is far from satisfiable}
Take $G$ to be the disjoint union of a perfectly satisfiable game and a game of value near $0$, each drawn with probability $\frac12$.  A function $h$ supported on the satisfiable half has $\|P_Gh\|=\|P_{T_G}h\|$, so $\lamp(G)=1$ although $\val(G)\approx\frac12$.
Theorem~\ref{thm:product} is then vacuous: the parameter $\lamp(G)$ cannot see that half of the game is hopeless, because one global denominator ignores how the game is connected.  The remedy of \citet{DinurSteurer14} for general games is to replace $T_G$ by one trivial game per question (their relaxation $\val_+$);
for the games that matter for Label Cover there is a simpler remedy, the subject of the next section: make the game expand.
\end{pitfall}


%% ===========================================================================
\section{Expanding games: the proof}
\label{sec:expander}
%% ===========================================================================

\subsection{The Markov chain of a game}

In the collision game of Lemma~\ref{lem:collgame} the referee draws two edges $e,e'$ at a common Bob-question.  Let $J_G(u,u')=\Pr[u_e=u,\ u_{e'}=u']$ in that experiment.  It is a symmetric matrix on $U$ whose row sums are $\mu_U(u)$, because each of $u_e,u_{e'}$ alone has law $\mu_U$.
Let $A_G(u,u')=J_G(u,u')/\mu_U(u)$ be the corresponding random walk: from $u$, move to a random Bob-question $w\sim\mu(\cdot\mid u)$ and come back to a random neighbour $u'\sim\mu(\cdot\mid w)$.  Write $\langle x,y\rangle_U=\E_{u\sim\mu_U}x(u)y(u)$ for functions on $U$ and $\mathrm{Var}(x)=\langle x,x\rangle_U-\langle x,\mathbf 1\rangle_U^2$.
Then $\langle x,A_Gy\rangle_U=\sum_{u,u'}J_G(u,u')x(u)y(u')$, so $A_G$ is self-adjoint with respect to $\mu_U$, and
\[
  \langle x,A_Gx\rangle_U=\E_w\Big(\E_{e\sim\mu(\cdot\mid w)}x(u_e)\Big)^2\ \ge\ 0 .
\]
Hence $A_G$ has eigenvalues $1=\lambda_1\ge\lambda_2\ge\dots\ge0$, and the quantity
\[
  \gamma(G)\;=\;1-\lambda_2(A_G)\ \in[0,1]
\]
is the \defn{spectral gap} of the game; $G$ is \defn{$\gamma$-expanding} if $\gamma(G)\ge\gamma$.  Since $\langle x,(I-A_G)x\rangle_U=\E_w\mathrm{Var}_{e\mid w}\,x(u_e)$ (the first term is $\E_ux(u)^2$ and the second is the square of the mean in each neighbourhood), the spectral gap is the best constant in a Poincar\'e inequality:
\begin{equation}\label{eq:poincare}
  \gamma(G)\cdot\mathrm{Var}_{u\sim\mu_U}(x)\ \le\ \E_w\,\mathrm{Var}_{e\sim\mu(\cdot\mid w)}\,x(u_e)\qquad\text{for all }x\colon U\to\R.
\end{equation}
In words: a function that varies a lot over $U$ must vary inside many neighbourhoods; a disconnected game has $\gamma=0$.

\begin{example}\label{ex:gaps}
(a) The CHSH game is \emph{free}: $\mu=\mu_U\otimes\mu_W$.  Then $J_G(u,u')=\mu_U(u)\mu_U(u')$, $A_G$ maps every function to its mean, and $\gamma=1$.  The same holds for every free game (Exercise~\ref{ex:free}).
(b) In the clause--variable game of Figure~\ref{fig:cv} (three clauses, four variables, nine equally likely edges) one finds $J(C_i,C_i)=\frac4{27}$ and $J(C_i,C_j)=\frac5{54}$ for $i\ne j$, so $A_G=\frac16I+\frac5{18}\mathbf 1\mathbf 1^\top$ has second eigenvalue $\frac16$ and $\gamma=\frac56$ (Exercise~\ref{ex:fig-gap}).
(c) For the odd cycle game on $\ell$ vertices, $\gamma=1-\cos^4(\pi/\ell)\approx2\pi^2/\ell^2$ (Exercise~\ref{ex:cycle-gap}); the triangle has $\gamma=\frac{15}{16}$.
\end{example}

\subsection{Expansion makes the product theorem effective}

\begin{theorem}[$\lamp$ is close to the collision value on expanding games; Dinur--Steurer]\label{thm:expander}
Let $G$ be a $\gamma$-expanding projection game and let $\eps>0$ satisfy $\lamp(G)>1-\eps$.  Then $\ \|G\|^2\ge1-\dfrac{54\,\eps}{\gamma}.$
\end{theorem}
\begin{proof}
Pick $h\ge0$ with $\|P_Gh\|>(1-\eps)\|P_{T_G}h\|$, and scale it so that $\|P_{T_G}h\|=1$.  Write $m(u)=\sum_ah(u,a)$; then $\E_um(u)^2=1$ and $\|P_Gh\|>1-\eps$.

\emph{Step 1: we may assume $h$ is concentrated on one label per question.}  Choose independently for every $u$ with $m(u)>0$ a label $a_u$ with probability $h(u,a)/m(u)$, and set $h'(u,a_u)=m(u)$, $h'(u,a)=0$ otherwise.  Then $\E h'=h$ and $m$ is unchanged, so $\|P_{T_G}h'\|=1$; by convexity of the norm, $\E\|P_Gh'\|\ge\|P_G\E h'\|=\|P_Gh\|$, so some outcome $h'$ has $\|P_Gh'\|\ge\|P_Gh\|$.  (This is where non-negativity of $h$ is used.)
Henceforth $h=m\cdot f_\ell$ for a labeling $\ell\colon U\to\Sigma_U$.

\emph{Step 2: the quadratic form.}  For $u,u'$ with $J(u,u')>0$ let $Q(u,u')\in[0,1]$ be the probability, in the experiment of Lemma~\ref{lem:collgame}, that $\pi_e(\ell(u))=\pi_{e'}(\ell(u'))$ given $u_e=u$, $u_{e'}=u'$.  The computation of Lemma~\ref{lem:collgame}, now with the weights $m$, gives
\begin{equation}\label{eq:coll-quad}
  \|P_Gh\|^2=\sum_{u,u'}J(u,u')\,m(u)m(u')\,Q(u,u'),\qquad
  \|G\|^2\ \ge\ \|P_Gf_\ell\|^2=\sum_{u,u'}J(u,u')\,Q(u,u').
\end{equation}
Write $\E_J$ for the average over $(u,u')\sim J$; so $\|G\|^2\ge\E_JQ$ and $\E_J[mm'\,Q]>(1-\eps)^2$, where $mm'$ abbreviates $m(u)m(u')$.

\emph{Step 3: $m$ is nearly constant.}  Since $Q\le1$ and $\E_J[mm']=\langle m,A_Gm\rangle_U$,
\[
  \langle m,(I-A_G)m\rangle_U=1-\E_J[mm']\le1-\E_J[mm'Q]<1-(1-\eps)^2\le2\eps .
\]
By~\eqref{eq:poincare}, and since $\mathrm{Var}_{e\mid w}x(u_e)$ averaged over $w$ is exactly $\langle x,(I-A_G)x\rangle_U$,
$\ \mathrm{Var}(m)\le2\eps/\gamma$.  Put $\bar m=\E_um(u)$; then $\bar m^2=1-\mathrm{Var}(m)$.  If $\eps/\gamma>\frac16$ the claim is trivial (its right side is negative), so assume $\eps/\gamma\le\frac16$; then $\mathrm{Var}(m)\le\frac13$ and $\bar m\ge\frac23$.

\emph{Step 4: the conclusion.}  For $s,t\ge0$ and $c\ge\frac23$,
\begin{equation}\label{eq:pointwise}
  1\ \le\ 9st+9(s-c)^2+9(t-c)^2 .
\end{equation}
Indeed, if $s<\frac13$ then $(s-c)^2\ge(\frac23-\frac13)^2=\frac19$; likewise for $t$; and if both are at least $\frac13$ then $9st\ge1$.
Apply~\eqref{eq:pointwise} with $c=\bar m$, $s=m(u)$, $t=m(u')$, multiply by $1-Q(u,u')\in[0,1]$ and average over $J$:
\[
  1-\E_JQ\ \le\ 9\,\E_J\big[(1-Q)mm'\big]+9\,\E_J\big[(1-Q)(m(u)-\bar m)^2\big]+9\,\E_J\big[(1-Q)(m(u')-\bar m)^2\big].
\]
Both marginals of $J$ equal $\mu_U$ and $0\le1-Q\le1$, so each of the last two terms is at most $9\,\mathrm{Var}(m)\le18\eps/\gamma$.  The first is $9(\E_J[mm']-\E_J[mm'Q])\le9(1-(1-\eps)^2)\le18\eps$, using $\E_J[mm']=\langle m,A_Gm\rangle_U\le\langle m,m\rangle_U=1$.
Hence $1-\|G\|^2\le1-\E_JQ\le18\eps+36\eps/\gamma\le54\eps/\gamma$, as $\gamma\le1$.
\end{proof}

\begin{corollary}[Parallel repetition for expanding projection games]\label{cor:expander}
Let $G$ be a $\gamma$-expanding projection game with $\val(G)\le1-\eta$.  Then $\lamp(G)\le1-\gamma\eta/54$, and for all $k\ge1$
\[
  \val(G^{\otimes k})\ \le\ \|G^{\otimes k}\|\ \le\ \Big(1-\frac{\gamma\eta}{54}\Big)^k .
\]
\end{corollary}
\begin{proof}
Let $\eps>0$ with $\lamp(G)>1-\eps$.  By Theorem~\ref{thm:expander} and Lemma~\ref{lem:val-coll}, $1-54\eps/\gamma\le\|G\|^2\le\val(G)\le1-\eta$, so $\eps\ge\gamma\eta/54$.  As this holds for every $\eps>1-\lamp(G)$, we get $1-\lamp(G)\ge\gamma\eta/54$.  The rest is Corollary~\ref{cor:lamp}.
\end{proof}

\noindent
The rate is linear in $\eta$, not quadratic as in the general bounds of Section~\ref{sec:general}; the price is the factor $\gamma$.  Numerical experiments for these notes (876 random projection games with $\gamma>0.001$ and $\lamp<1$, half of them planted to be nearly satisfiable) never produced a value of $(1-\|G\|^2)\gamma/\eps$ above $3.7$, where $\eps=1-\lamp(G)$; the proved bound is $54$.  The constant is certainly not tight, and we have not tried to optimize it.

\begin{pitfall}{what the expansion hypothesis is doing}
Steps~3--4 are the only places where $\gamma$ enters.  A function $h$ with $\|P_Gh\|\approx\|P_{T_G}h\|$ could, a priori, be a large multiple of a perfect labeling on a tiny part of $U$ and nearly vanish elsewhere; then the weights $m$ would be far from constant and $\|G\|$ could be small even though $\lamp(G)\approx1$.
That is exactly the disconnected game of the pitfall above.  Expansion forbids it: a function with small Dirichlet form $\langle m,(I-A_G)m\rangle_U$ is close to a constant.
\end{pitfall}


%% ===========================================================================
\section{The Label Cover theorem, with explicit parameters}
\label{sec:lc}
%% ===========================================================================

The hypothesis of Corollary~\ref{cor:expander} is easy to arrange: we may mix in a harmless piece of game that connects everything.

\begin{lemma}[Making a game expand]\label{lem:mix}
Let $G$ be a projection game, $0\in\Sigma_W$ a fixed label.  Let $G'$ be the game with one new Bob-question $w_0$ whose edge distribution is
\[
  \mu'=\tfrac12\mu\ +\ \tfrac12\sum_{u\in U}\mu_U(u)\,\delta_{(u,w_0)},
\]
where the edge $(u,w_0)$ carries the constant map $\pi(a)=0$.  Then (i)~$\val(G')=\frac12+\frac12\val(G)$, and (ii)~$\gamma(G')\ge\frac12$.  The question sets $U$, the label sets, and the maps of the old edges are unchanged.
\end{lemma}
\begin{proof}
(i) Any strategy wins the new edges with probability at most $\frac12$, their total weight, and the old edges with probability at most $\frac12\val(G)$.  Both bounds are attained together: Bob answers $0$ at $w_0$ (which wins every new edge, whatever Alice does) and plays an optimal strategy of $G$ elsewhere.
(ii) The marginal of $\mu'$ on $U$ is $\mu_U$ again.  In the collision experiment of $G'$, with probability $\frac12$ the Bob-question is $w_0$ and $u,u'$ are independent copies of $\mu_U$; otherwise it is as in $G$.  Thus $A_{G'}=\frac12A_G+\frac12\mathbf 1\mu_U^\top$.  On functions orthogonal to the constants $A_{G'}$ acts as $\frac12A_G$, so $\lambda_2(A_{G'})=\frac12\lambda_2(A_G)\le\frac12$.
\end{proof}

\begin{remark}[Uniformity]
If $G$ is uniform on $|E|$ edges, then $\mu_U(u)=\deg(u)/|E|$ and $G'$ is the \emph{uniform} game on $2|E|$ edges obtained by adding $\deg(u)$ copies of $(u,w_0)$ for every $u$.  So no weights are needed, and Remark~\ref{rem:multiplicity} applies to $G'^{\otimes k}$.
\end{remark}

\begin{theorem}[Label Cover theorem, with constants]\label{thm:lc-explicit}
Suppose it is $\NP$-hard to distinguish projection games (with label sets of sizes $7$ and $2$) of value $1$ from those of value at most $1-\eta_0$; by Theorem~\ref{thm:pcp2} and Proposition~\ref{prop:cv} this holds with $\eta_0=\eps_0/3$.  Then for every $\delta\in(0,1)$ and
\[
  k\ =\ \Big\lceil\frac{216}{\eta_0}\ln\frac1\delta\Big\rceil,
\]
the map $G\mapsto(G')^{\otimes k}$ is a polynomial-time reduction (for $\delta$ fixed) to distinguishing projection games with label sets of sizes $7^k$ and $2^k$ of value $1$ from value at most~$\delta$.  In particular $7^k,2^k=\poly(1/\delta)$.
\end{theorem}
\begin{proof}
If $\val(G)=1$ then $\val(G')=1$ and $\val((G')^{\otimes k})=1$ by Lemma~\ref{lem:rep-basic}.  If $\val(G)\le1-\eta_0$ then $\val(G')\le1-\eta_0/2$ by Lemma~\ref{lem:mix}(i), and $G'$ is $\frac12$-expanding by (ii).  Corollary~\ref{cor:expander} with $\gamma=\frac12$ and $\eta=\eta_0/2$ gives
$\val((G')^{\otimes k})\le(1-\eta_0/216)^k\le e^{-k\eta_0/216}\le\delta$.  The product has $(2|E|)^k$ edges, polynomial for fixed $\delta$.
\end{proof}

\noindent
This is Theorem~\ref{thm:lc} again, now with a proof that is complete except for the PCP theorem, with $k=O(\log(1/\delta))$ and exponent $\poly(1/\delta)$ of order $1/\eta_0$.
The proof in Lecture~\ref{lec:why} cited Rao's bound, which needs no expansion but has rate $\eta^2$ instead of $\eta$.  Here the number of repetitions is linear in $1/\eta_0$ rather than quadratic, at the price of the expansion step of Lemma~\ref{lem:mix} and of a worse absolute constant; and nothing beyond linear algebra is used.

\begin{pitfall}{the Label Cover theorem does not give Unique Games}
Taking a product of projection maps preserves their fibre sizes: if each $\pi_e$ is $d$-to-$1$ then each $\pi_{e_1}\times\dots\times\pi_{e_k}$ is $d^k$-to-$1$, and a product of bijections is a bijection.
The games produced by Theorem~\ref{thm:lc-explicit} are therefore $d$-to-$1$ with $d$ growing with $k$ (for the clause--variable game, the fibres have $3$ or $4$ elements; the mixed-in edges are constant maps).
Not one of them is a Unique Game, and repetition cannot make them one.  We return to this in Section~\ref{sec:ugc-rep}.
\end{pitfall}


%% ===========================================================================
\section{What is true in general, and what is not}
\label{sec:general}
%% ===========================================================================

Expansion was a convenience.  The theorems below need no such hypothesis; they are all stated for $\val(G)\le1-\eta$ unless noted.

\begin{table}[t]
\centering\small
\renewcommand{\arraystretch}{1.3}
\begin{tabular}{@{}l>{\raggedright\arraybackslash}p{3.6cm}>{\raggedright\arraybackslash}p{4.9cm}@{}}
\toprule
result & games & bound on $\val(G^{\otimes k})$\\
\midrule
\citet{Raz98}; \citet{Holenstein09} & all games, answer sets of size $\le\sigma$ & $(1-\eta^{c})^{\Omega(k/\log\sigma)}$, with $c=32$, resp.\ $c=3$ (Theorem~\ref{thm:raz})\\
\citet{Rao11} & projection games & $(1-\Omega(\eta^2))^k$\\
\citet{DinurSteurer14} & projection games with $\val(G)\le\rho$, any $\rho\in(0,1)$ & $\big(\frac{2\sqrt\rho}{1+\rho}\big)^{k/2}$; this is $\le(1-\eta^2/16)^k$ for $\rho=1-\eta$ and $\le(4\rho)^{k/4}$ for small $\rho$\\
Corollary~\ref{cor:expander} & $\gamma$-expanding projection games & $(1-\gamma\eta/54)^k$\\
\bottomrule
\end{tabular}
\caption{Parallel repetition theorems.  The first three are quoted; the last was proved in this lecture.}
\label{tab:rep-thms}
\end{table}

\noindent
The new feature of the Dinur--Steurer bound is that $\rho$ may be as small as one likes; before it, it was not known that repetition decreases a value that is already sub-constant \citep[Section~1]{DinurSteurer14}.  Their proof is the general version of the argument of this lecture, with the relaxation $\val_+$ in place of $\lamp$.

\paragraph{Strong repetition.}  The bounds of Table~\ref{tab:rep-thms} for projection games decay like $(1-\Omega(\eta^2))^k$.  \defn{Strong parallel repetition} is the hope that $(1-\Omega(\eta))^k$ holds, as it would if the na\"ive guess were roughly right.
It is \emph{false} in general, even for Unique Games on two labels: the odd cycle game of Lecture~\ref{lec:why} has $\eta=\frac1{2\ell}$, yet $G^{\otimes k}$ keeps value $1-O(\sqrt k/\ell)$ up to $k\approx\ell^2$ \citep{Raz11}, and this is essentially tight \citep{FKO07}.
It is true for restricted classes: free projection games \citep{BRRRS09} and projection games on expanders \citep{RR12}.  Corollary~\ref{cor:expander} has the same flavour, with a rate linear in $\eta$ but proportional to the gap $\gamma$; for free games $\gamma=1$ (Exercise~\ref{ex:free}).
The two facts are consistent: the odd cycle has $\gamma\approx2\pi^2/\ell^2$ (Exercise~\ref{ex:cycle-gap}), so Corollary~\ref{cor:expander} only promises decay after $k\gtrsim\ell^3$ repetitions.
As \citet{DinurSteurer14} note, a strong repetition bound for every Unique Game ``would imply a reduction from max cut to unique games''; Raz's counterexample rules out this route.

\begin{bigpicture}{The shape of the information-theoretic proof (a sketch; none of this is proved here)}
The theorems of Raz and Holenstein are proved by contradiction, one coordinate at a time.  Fix a strategy for $G^{\otimes k}$, and for $C\subseteq[k]$ let $W_C$ be the event that all coordinates in $C$ are won.
The central claim is: \emph{if $\Pr[W_C]$ is not too small (say $\ge2^{-\beta k}$), there is a coordinate $i\notin C$ with $\Pr[W_i\mid W_C]\le\val(G)+o(1)$}.  Iterating, $\Pr[W_{C\cup\{i\}}]\le(\val(G)+o(1))\Pr[W_C]$ until the probability falls below $2^{-\beta k}$, which gives exponential decay.
Why the claim holds: conditioning on an event of probability $2^{-\beta k}$ can distort the $k$ pairs of questions by at most $\beta k$ bits of information in total, so for most coordinates $i$ the pair $(x_i,y_i)$ is still almost distributed as in $G$.
The two provers can then \emph{simulate} a strategy for the single game $G$ at coordinate $i$: they embed the real questions at coordinate $i$, and use shared randomness (correlated sampling) to fill in the other coordinates, together with the answers, from the conditional distribution given $W_C$, each looking only at its own question.
The simulated strategy wins $G$ with probability about $\Pr[W_i\mid W_C]$, hence this is at most $\val(G)+o(1)$.  The dependence on the size $\sigma$ of the answer sets enters through the same information accounting.
\end{bigpicture}

\noindent
Compare the two proofs.  Raz and Holenstein follow a \emph{strategy}: they transform a strategy for the product into one for a single copy.  The proof of this lecture never mentions a strategy of the product; it bounds the \emph{norm} of an operator and uses a spectral property of a graph.
The first is more general; the second shorter, and it yields the explicit rate $\gamma\eta$ that makes the role of the graph visible.


%% ===========================================================================
\section{What repetition cannot do for Unique Games}
\label{sec:ugc-rep}
%% ===========================================================================

Repetition turned the PCP theorem into Label Cover.  The Unique Games Conjecture asks for a Label Cover \emph{with bijections}, completeness $1-\zeta$ and soundness $\delta$, for every $\zeta,\delta>0$.  It is natural to try the same recipe, and it is instructive to see exactly where it breaks.

\begin{bigpicture}{Three reasons why repetition does not prove the Unique Games Conjecture}
\textbf{1.\ The type of the constraints is preserved.}  A product of $d$-to-$1$ maps is $d^k$-to-$1$ and a product of bijections is a bijection (Pitfall in Section~\ref{sec:lc}).  The PCP theorem gives projections that are not bijections (fibres of size $3$ and $4$ for $3$-CNF), and repetition cannot repair that.

\smallskip
\textbf{2.\ Repetition never improves completeness.}  By Lemma~\ref{lem:rep-basic}, $\val(G^{\otimes k})\le\val(G)$ for every instance, so whatever completeness the base instances have, repetition cannot raise it.  Perfect completeness is not available as a base, since Unique Games of value $1$ can be recognized in polynomial time (Proposition~\ref{prop:ug-easy}).
The UGC asks for completeness $1-\zeta$ for \emph{every} $\zeta>0$; a repetition-based proof would have to start from a hardness result for Unique Games with the same completeness and a weaker, constant soundness, that is, from a result of the same kind as the conjecture.

\smallskip
\textbf{3.\ Soundness does not decay at the na\"ive rate.}  Even from such a base problem, $(1-\Omega(\eta))^k$ is false for Unique Games \citep{Raz11,FKO07}: the odd cycle games keep value $1-O(\sqrt k\,\eta)$ for $k$ up to about $1/\eta^2$, so in the worst case about $1/\eta^2$ repetitions are needed before the value drops appreciably.
Starting from completeness $1-\zeta$ and value at most $1-\eta$ on the no-instances, $k$ repetitions give completeness at most $1-\zeta$ and, by Rao's bound, soundness at most $e^{-\Omega(k\eta^2)}$.

\smallskip
These are not proofs that the conjecture is hard to prove; they say that the proof, if it exists, must use a different idea.  Lectures~8--12 use the conjecture as a hypothesis, and Lectures~13--15 prove its $2$-to-$2$ relative, which is \emph{not} a statement about Unique Games but comes close.
\end{bigpicture}


%% ===========================================================================
\section*{Exercises}
\addcontentsline{toc}{section}{Exercises}
%% ===========================================================================

\begin{exercise}[\stars{$\star\star$}]\label{ex:rep-prod}
\emph{Products and an experiment.}  (a)~Prove $\val(G)\val(H)\le\val(G\otimes H)\le\min\{\val(G),\val(H)\}$.  (b)~Prove $\|G\otimes H\|\ge\|G\|\,\|H\|$ by exhibiting a product assignment.
(c)~Write a program that computes $\val(G)$, $\|G\|$ and $\lamp(G)$ for small random projection games, and test Theorem~\ref{thm:product}, and compare $\|G\otimes H\|$ with $\|G\|\,\|H\|$.
\emph{Hints for (c):} for a labeling $\ell$ put $c_\ell(w,b,u)=\sum_{e=(u,w)\colon\pi_e(\ell(u))=b}\mu(e)$; then $\|P_Gf_\ell\|^2=\sum_{w,b}\big(\sum_uc_\ell(w,b,u)\big)^2/\mu_W(w)$, and $\lamp(G)^2$ is the maximum over $\ell$ of the top eigenvalue of $D^{-1/2}M_\ell D^{-1/2}$, with $D=\mathrm{diag}(\mu_U)$ and $M_\ell(u,u')=\sum_{w,b}c_\ell(w,b,u)c_\ell(w,b,u')/\mu_W(w)$.
\emph{Warning:} most small random games are satisfiable, hence have $\|G\|=1$, and all the inequalities are trivial for them; filter for $\|G\|<1$ (dense graphs with random projection maps work).  Does $\|G\otimes H\|<\|G\|\,\|H\|$ ever occur in your experiments?  These notes make no claim either way.
\end{exercise}

\begin{exercise}[\stars{$\star$}]\label{ex:fig-gap}
For the clause--variable game of the formula of Figure~\ref{fig:cv} compute the matrices $J$ and $A_G$ and verify $\gamma(G)=\frac56$.  Apply Lemma~\ref{lem:mix} and verify that $\gamma(G')=\frac{11}{12}$ (so the bound $\frac12$ of the lemma is far from tight here).
\end{exercise}

\begin{exercise}[\stars{$\star\star$}]\label{ex:cycle-gap}
Let $G$ be the odd cycle game on $\ell$ vertices.  (a)~Show that $A_G=K^2$ for $K=\frac12I+\frac14(S+S^{-1})$, where $S$ is the cyclic shift, and that $A_G$ has eigenvalues $\cos^4(\pi j/\ell)$, $0\le j<\ell$.  Conclude $\gamma(G)=1-\cos^4(\pi/\ell)$, which is $\frac{15}{16}$ for the triangle.
(b)~With $\eta=\frac1{2\ell}$, Corollary~\ref{cor:expander} gives $\val(G^{\otimes k})\le(1-\gamma\eta/54)^k$.  Show that this is at least $\frac12$ for $k\le\ell^3$ and $\ell$ large, so that it does not contradict Raz's example (which keeps the value above $\frac12$ for $k\le c\ell^2$).  This is a consistency check, not a claim that the corollary is tight.
\end{exercise}

\begin{exercise}[\stars{$\star\star$}]\label{ex:disjoint}
Let $G_1$ be a satisfiable projection game, $G_2$ any projection game with $\val(G_2)<1$ on question sets disjoint from those of $G_1$, and let $G$ draw an edge of $G_1$ or of $G_2$ with probability $\frac12$ each.  Show that $\lamp(G)=1$ and $\gamma(G)=0$, while $\|G\|^2\le\val(G)\le\frac12+\frac12\val(G_2)<1$.
Conclude that no statement of the form of Theorem~\ref{thm:expander} can hold without a hypothesis on $\gamma$.  Which step of its proof fails?
\end{exercise}

\begin{exercise}[\stars{$\star$}]\label{ex:ug-mix}
Let $G$ be a Unique Game.  (a)~Show that the game $G'$ of Lemma~\ref{lem:mix} is not a Unique Game, and explain why the proof of Theorem~\ref{thm:lc-explicit} therefore gives no information about Unique Games.
(b)~Suppose the new edges $(u,w_0)$ carried bijections $\sigma_u$ instead of the constant map.  Show that Lemma~\ref{lem:mix}(i) then fails: find a satisfiable $G$ for which $\val(G')<1$.  \emph{Hint: a path $u_1$---$w$---$u_2$ with identity maps, and $\sigma_{u_1}=\mathrm{id}$, $\sigma_{u_2}=$ the transposition of the two labels.}  Why does this destroy the argument?
\end{exercise}

\begin{exercise}[\stars{$\star\star$}]\label{ex:free}
A game is \emph{free} if $\mu=\mu_U\otimes\mu_W$.  Show that $J_G=\mu_U\otimes\mu_U$ and $\gamma(G)=1$ for every free game.  Deduce $\val(G^{\otimes k})\le(1-\eta/54)^k$ for every free projection game of value at most $1-\eta$.  Check on CHSH that $\val(G^{\otimes k})>\val(G)^k$ for $k=2,3$, and that it nevertheless respects the bound.
\end{exercise}

\begin{exercise}[\stars{$\star\star\star$}]\label{ex:reading-ds}
\emph{Reading.}  Read Section~4 of \citet{DinurSteurer14}.  Identify the place where our proof of Theorem~\ref{thm:expander} uses expansion, and explain how their ``partial assignments'' together with correlated sampling replace it for games with $\gamma=0$.
Then compare the dependence on $\eta$ in Corollary~\ref{cor:expander} (linear) with that of Corollary~1.2 of their paper (quadratic) and say what accounts for the difference.
\end{exercise}


%% ===========================================================================
\section*{Notes and further reading}
\addcontentsline{toc}{section}{Notes and further reading}
%% ===========================================================================

\paragraph{Where this lecture comes from.}
Sections~\ref{sec:operator}--\ref{sec:lc} follow Sections~2 and~3 and Claim~A.1 of \citet{DinurSteurer14} closely: the collision value, Claims~2.3 and~2.4 (our Lemmas~\ref{lem:val-coll} and~\ref{lem:coll-mono}), the parameter $\lamp$ and Theorems~3.2 and~3.4 (our Theorems~\ref{thm:product} and~\ref{thm:expander}), and the expansion trick.
The intuition that one should compare $G\otimes H$ with $T\otimes H$ to ``cancel $H$'' is theirs.  Section~\ref{sec:general} and Table~\ref{tab:rep-thms} summarize the literature as described in the introduction of that paper.

\paragraph{What we changed.}
(i)~\emph{Orientation.}  Dinur and Steurer let the player with the large alphabet be Bob; we let it be Alice, as in Lectures~\ref{lec:why} and~\ref{lec:games}, so the roles of ``$u$'' and ``$v$'' are swapped throughout and the new vertex of the expansion trick sits on Bob's side.
(ii)~\emph{Totality.}  Their projections may be partial; ours are total.  This removes the step ``increase some entries of $P_Gf$ to turn it into an assignment'' in the proof of $\|G\|^2\le\val(G)$, and it is why Feige's game (Example~\ref{ex:feige}) is outside our class.
(iii)~\emph{Constants.}  They write $\|G\|>1-O(\eps/\gamma)$; the proof ends in a bound on $\|G\|^2$, which is the form we state and the one we need, with the explicit constant $54$.
(iv)~\emph{Presentation.}  We avoid $\val_+$ altogether, since the expanding case does not need it.

\paragraph{On attributions.}
Raz's theorem appeared in preliminary form at STOC~1995 and in full as \citet{Raz98}; Holenstein's simplification and Rao's projection-game bound are cited as in Lecture~\ref{lec:why}.
Feige's non-interactive agreement game is from \citet{Feige91}, as presented in Appendix~B of \citet{DinurSteurer14}, where $\val(G)=\val(G^{\otimes2})=\frac12$ is stated; we wrote out the short proofs and confirmed both values by exhaustive search.
The CHSH game is the XOR game of Clauser--Horne--Shimony--Holt; its classical values $\frac34,\frac58,\frac{31}{64}$ for $k=1,2,3$ were computed for these notes by exhaustive search, with $k=3$ checked by two separate programs.
The odd cycle game and its repetition are due to \citet{FKO07} and \citet{Raz11}.

\paragraph{What is not covered.}
The proof of the repetition theorem for non-expanding games (the relaxation $\val_+$, the Cheeger-type inequality, correlated sampling), the information-theoretic proofs of Raz and Holenstein, the tight examples for the dependence on the alphabet size, and the applications to Set Cover and to sub-constant soundness (see \citealp{DinurSteurer14}).
Lecture~6 turns to the other classical way of amplifying a gap, due to \citet{Dinur07}; unlike repetition it keeps the alphabet small.
